\section{What the Loader Does Not Fix}
\label{sec:ch10-what-the-loader-does-not-fix}

\index{projection!the entity route has none|(}%
Look at that batched statement again, next to what the node resolver in the
same file emits for a request asking for the same one field:

\begin{minted}{text}
-- the entity route, asked for title
SELECT "s"."Id", "s"."Abstract", "s"."DurationMinutes",
       "s"."SpeakerId", "s"."StartsAt", "s"."Title"
FROM "Sessions" AS "s"
WHERE "s"."Id" IN (@ids1, @ids2, @ids3)

-- the node route, asked for title
SELECT "s"."Title"
FROM "Sessions" AS "s"
WHERE "s"."Id" = @id
LIMIT 1
\end{minted}

\index{projection!six columns against one}%
Six columns against one. The loader fixed the number of statements and did
nothing at all about the width of them, and the entity route is the route the
router uses for every cross-seam field in the graph. So this is the more common
of the two paths taking the worse of the two shapes.

\index{projection!chapter 6 said it was impossible here}%
Chapter~\ref{ch:first-subgraph} said this was not fixable, and printed a comment
in \code{SessionType.cs} saying a \code{QueryContext} could not be injected into
a reference resolver. That is wrong. It can, and the comment in the file has
changed accordingly, which is the version section~\ref{sec:ch10-the-loader-behind-the-resolver}
prints.

\subsection{The Projection That Works}

\index{QueryContext@\code{QueryContext<T>}!does inject into a reference resolver}%
Take the reference resolver as chapter~\ref{ch:first-subgraph} left it, add a
\code{QueryContext<Session>} parameter, and put chapter~\ref{ch:data-without-n-plus-1}'s
\code{.With(query)} back on the query. It builds, it starts, and it answers,
and the statement it sends selects one column rather than six: the same
\code{SELECT "s"."Title"} the node route beside it emits. Ask for two fields and
it widens to two. Nothing here is a state this book ships, so I am describing
these runs rather than printing them; what each one issued is recorded in the
chapter's research note.

\index{QueryContext@\code{QueryContext<T>}!the selector it carries}%
The selector Hot Chocolate handed that resolver is exactly right. It constructs
a \code{Session} taking \code{Title} from the row and defaulting the other five
properties, which is what a projection for a query asking only for the title
should look like.

\index{ISelection@\code{ISelection}!at the entity route is \code{\_entities}}%
Which is more surprising than it looks, because the selection the resolver is
standing in is not the one that selector describes. Ask the resolver what field
it is resolving and it says \code{\_entities}, declared on \code{Query},
returning \code{[\_Entity]!}, whose named type is a union. The engine built the
selector from the inline fragment inside that field while telling the resolver
it is resolving the root. Anything that reads the selection by hand gets the
union; the injected \code{QueryContext} gets what you wanted.

\subsection{And Why the Book Does Not Ship It}

\index{projection!the selector defaults the key}%
Read that selector once more, at its first argument. Nobody asked for \code{id},
so the projection fills it with \code{default(Int32)}, and the object that comes
back has \code{Id} zero. For a resolver returning one row that costs nothing:
the row is used for the fields that were asked for and the key was only ever the
thing in the \code{WHERE} clause.

\index{DataLoader!a dictionary keyed on a column the projection dropped}%
For a loader it is fatal. Put the projection inside the loader and the query
that goes out is everything you wanted: one statement, one column, and the
three keys in an \code{IN} list. Then look at what comes back. Three
\code{Speaker} records with the right names on them and \code{Id} zero on all
three, and a dictionary keyed on \code{Id} being built out of them, which ends
where you would expect: an argument exception saying an item with that key has
already been added, naming the key as zero.

\index{DataLoader!keyed on a column the projection dropped}%
The service answers \code{Unexpected Execution Error} for every representation.
Chapter~\ref{ch:second-third-subgraph} got that same message out of a stub that
advertises a key it cannot resolve, and out of a service missing a serializer
registration. Here it is a dictionary with a duplicate key in it. Three
unrelated faults, one sentence, and it never says which.

\index{projection!correct only if the client asked for the key}%
Here is the part that decides it. Add \code{id} to the client's selection set
and the same code works perfectly. The selector picks the key up because
somebody asked for it, the statement gains a column, the dictionary is keyed 1,
2 and 3, and the answer is correct. A subgraph that is right when the caller
asks for the key and throws when it does not is not a subgraph I am willing to
print, and it is not a failure anybody would catch in review, because the
request you test with is the one that has the id in it.

\index{records!positional records block the API built for this}%
GreenDonut has an API for exactly this problem, a \code{Select} overload whose
first argument is a key selector precisely so that the key survives the
projection. There are two of them, one taking a builder and one taking an
expression, and neither runs here: both reach a rewrite that calls
\code{Expression.New} on the entity type, and both raise an argument exception
saying \code{Speaker} has no default constructor.

\index{records!the projection wall, twice}%
Which is chapter~\ref{ch:data-without-n-plus-1}'s wall again. That chapter could
not use \code{[UseProjection]}, because it builds its projected row the same
way and a positional record has no parameterless constructor. It is the same
rewrite failing for the same reason, which closes the whole selector-builder
family to a domain modeled in records, and this book's domain has been records
since chapter~\ref{ch:hot-chocolate-zero}. There is a third route, merging a
key-preserving expression into the builder by hand, and it is worse than a
failure: it raises nothing and drops the expression, and I cannot tell you why.

\index{projection!batch or narrow, not both}%
So the entity route can be batched, always, and it can be projected, at the cost
of the batch. It cannot be both, and given the choice I take the batch: the
number of statements grows with somebody else's page size and the number of
columns does not. That is what the comment in \code{SessionDataLoader.cs} is
recording, and it is there rather than in this paragraph because a reader typing
the file out needs the reason at the line.

\index{nullability!a wide select is not free forever}%
One honest caveat about the trade. Six narrow columns on a SQLite table is not a
cost anybody can measure, and I am not going to pretend it is. It becomes one
when a row is wide, when a column is a blob nobody asked for, or when the
projection would have let the database answer from an index instead of the
table. If that is your shape, the entity route is where to look, and the answer
this book could not find is still worth looking for.
\index{projection!the entity route has none|)}%
\index{N+1@N+1!at the entity route|)}%
